\documentclass[11pt]{article}
\usepackage{amsmath,amssymb,amsthm}
\usepackage[margin=1.1in]{geometry}
\newtheorem{theorem}{Theorem}
\newtheorem{lemma}[theorem]{Lemma}
\newtheorem{corollary}[theorem]{Corollary}
\newtheorem{proposition}[theorem]{Proposition}
\newtheorem{remark}[theorem]{Remark}
\newcommand{\Sym}{\operatorname{Sym}}
\newcommand{\Gm}{\Gamma_{\mathbb C}}
\newcommand{\Gr}{\Gamma_{\mathbb R}}
\title{Pole detection at the carrier:\\
the beyond-endoscopy dichotomy for transported representation data}
\author{Samuel Lavery\thanks{Machine-checked companion:
\texttt{RequestProject/PoleDetection.lean}, 4 named declarations, each at axiom
footprint $\{\texttt{propext},\texttt{Classical.choice},\texttt{Quot.sound}\}$ or a
subset, no \texttt{sorry}, no \texttt{axiom}, zero \texttt{sorryAx}.}}
\date{Draft --- \today}

\begin{document}
\maketitle

\begin{abstract}
Langlands' beyond-endoscopy principle detects functorial origin through poles:
$L(s,\pi,r)$ has one at $s=1$ exactly when $r$ contains the trivial
representation.  This note compiles that detection for transported data
$\rho=\oplus_i\Sym^{r_i}$ of the level-one seed at the mechanism's package level.
Unconditionally, a trivial constituent splits a $\zeta$-factor off the
datum bank, explicit in the $L$-readout on the joint half-plane.  The detection
theorem: \emph{a datum containing the trivial representation admits no strict theta
reflection} --- from the strict joint identity the mechanism yields an entire
completed transform, while the $\zeta$-constituent forces the simple pole at $s=1$;
analytic continuation across $\{\operatorname{Re}s>1\}$ and a ray limit against the
residue collide the two, unless the completed erased factor vanishes at the edge,
which is the named classical input.  Contrapositive: a datum satisfying the strict,
mass-free reflection has no trivial constituent.  The masses of the weak reflection
pair are the carrier avatar of the pole: $\zeta$-containing objects reflect only with
mass terms.
\end{abstract}

\section{The $\zeta$-factorization}

Fix the level-one eigenform seed with Hecke data, and a representation datum
$\rho=[r_1,\dots,r_m]$ with $0\in\rho$.  Both statements are unconditional:

\begin{lemma}[\texttt{repBank\_zeta\_factor},
\texttt{LSeries\_repBank\_zeta\_factor}; machine-checked]\label{lem:zeta}
The datum bank splits as
\[
c_\rho=\zeta\star c_{\rho\smallsetminus\{0\}},
\]
the rank-zero clock being the unit system
(\texttt{symrBank\_zero\_eq\_zeta}); on the joint half-plane
\[
L(\rho,s)=\zeta(s)\cdot L(\rho\smallsetminus\{0\},s),
\]
with $\zeta$ the Riemann zeta function of Mathlib
(\texttt{LSeries\_zeta\_eq\_riemannZeta}).
\end{lemma}

\section{The detection theorem}

\begin{theorem}[\texttt{trivial\_constituent\_detection}; machine-checked]
\label{thm:detect}
Let $0\in\rho$.  Suppose, at some general Deligne chart $(\delta s,\mu s)$ with
nonnegative shifts:
\begin{enumerate}
\item the \emph{strict} joint theta identity of the prescribed $\rho$-readouts
(the input of the mechanism, with zero masses);
\item a package $Q$ for the erased datum $\rho\smallsetminus\{0\}$ at the same chart
(entire $\Lambda_Q$, functional equation, chart identification); and
\item edge nonvanishing: $\Lambda_Q(1)\neq0$.
\end{enumerate}
Then \textbf{False}.
\end{theorem}

The proof is the collision of two compiled facts.  From (1),
\texttt{rep\_general\_package\_of\_theta} yields an \emph{entire} $\Lambda_P$ with
$\Lambda_P(s)=G(s)\,L(\rho,s)$ on the joint half-plane; by Lemma~\ref{lem:zeta} and
the chart of $Q$,
\[
\Lambda_P(s)=\zeta(s)\cdot\Lambda_Q(s)
\qquad(\operatorname{Re}s\ \text{large}),
\]
and both sides are analytic on the connected open half-plane
$\{\operatorname{Re}s>1\}$, so the identity persists there
(\texttt{AnalyticOnNhd.eqOn\_of\_preconnected\_of\_eventuallyEq}).  Along the real
ray $s=1+t$, $t\to0^+$:
\[
(s-1)\,\Lambda_P(s)\longrightarrow 0
\qquad\text{while}\qquad
(s-1)\,\zeta(s)\,\Lambda_Q(s)\longrightarrow \Lambda_Q(1)
\]
(Mathlib's \texttt{riemannZeta\_residue\_one} and continuity of the entire factors),
whence $\Lambda_Q(1)=0$, contradicting (3).

\begin{corollary}[contrapositive register]
A representation datum of the seed satisfying the strict theta reflection at an
admissible chart --- together with an erased-datum package nonvanishing at the edge
--- contains no trivial constituent.  Equivalently: the trivial constituent is
\emph{detected by the reflection type} --- $\zeta$-containing objects reflect only
with mass terms, never strictly.  The masses $f_0,g_0$ of the weak reflection pair
are the carrier-side avatar of the pole.
\end{corollary}

\section{Register}

\begin{remark}[register]
\textbf{Proved (machine-checked).}  The $\zeta$-factorization of the datum bank and
its $L$-readout; the threshold bound; the detection theorem exactly as stated, with
its three inputs typed in the signature.

\textbf{Consumed from the corpus.}  The mechanism at the general chart
(\texttt{mechanismG} via \texttt{rep\_general\_package\_of\_theta}), datum
self-duality, the compiled bank calculus.  From Mathlib: the residue of $\zeta$ at
$s=1$, differentiability of $\zeta$ off $s=1$, analytic-continuation uniqueness.

\textbf{Cited, not formalized.}  The classical discharge of input (3) --- edge
nonvanishing of the completed erased factor --- is the Jacquet--Shalika class of
nonvanishing theorems \cite{JS}; for symmetric-power factors of low rank it is known
unconditionally through the Kim--Shahidi lifts \cite{KimShahidi}.  Input (2) is the
mechanism's own output given the erased datum's identity.

\textbf{Not claimed.}  No trace formula, no orbital integrals, no general
automorphic $\pi$: the scope is the transported data of the level-one
$\mathrm{GL}(2)$ seed, and ``beyond endoscopy'' names the \emph{detection principle}
compiled here --- pole $\leftrightarrow$ trivial constituent $\leftrightarrow$
reflection type --- not Langlands' full program \cite{Langlands}.

\textbf{Falsifier.}  Exhibiting a datum with a trivial constituent that satisfies
the strict identity, together with a nonvanishing erased package, would make the
compiled theorem produce \textbf{False} --- an inconsistency of the corpus with
Mathlib.  Classically, the constant term of the theta function of any
$\zeta$-containing object is nonzero, which is why no such datum exists; the theorem
is the compiled form of exactly that obstruction.
\end{remark}

\begin{thebibliography}{9}
\bibitem{Langlands} R.~P.~Langlands,
\emph{Beyond endoscopy}, in: Contributions to automorphic forms, geometry, and
number theory, Johns Hopkins Univ.\ Press, 2004, 611--697.
\bibitem{JS} H.~Jacquet, J.~A.~Shalika,
\emph{A non-vanishing theorem for zeta functions of $\mathrm{GL}_n$},
Invent.\ Math.\ \textbf{38} (1976/77), 1--16.
\bibitem{KimShahidi} H.~Kim, F.~Shahidi,
\emph{Functorial products for $\mathrm{GL}_2\times\mathrm{GL}_3$ and the symmetric
cube for $\mathrm{GL}_2$}, Ann.\ of Math.\ (2) \textbf{155} (2002), 837--893.
\end{thebibliography}

\end{document}
